\section{Reflection and finite configuration algebra}\label{sec:configuration}

Let a finite labelled configuration be invariant under conjugation.
Suppose $|y_j|\le1$, and choose
$w_j=w((x_j-X)/h)$ with $0\le w\le1$ and compact support, $h>0$.
Only labels in that support matter. Define the horizontal moment
measures and their first useful masses by
\begin{equation}\label{eq:qmoments}
 q_\ell=\sum_jw_jy_j^\ell\delta_{x_j},\qquad
 N_w=q_0(\R),\qquad Q_w=q_2(\R).
\end{equation}
Conjugation cancels each odd moment at each horizontal coordinate,
not merely after integration:
$q_{2r+1}=0$. This elementary observation determines exactly how much
information a quadratic transverse difference can retain.

\begin{proposition}[Quadratic pair reduction]\label{prop:quadratic}
Let $\chi$ be bounded, even and nonnegative. Then
\begin{align}
 S_2(\chi)&=\sum_{j\ne k}
    w_jw_k\chi(x_j-x_k)(y_j-y_k)^2 \notag\\
 &=2\int(\chi*q_0)(x)\dd q_2(x)
 \le 2\norm{\chi}_\infty N_wQ_w .\label{eq:quadratic-reduction}
\end{align}
In particular $S_2(1)=2N_wQ_w$.
\end{proposition}
\begin{proof}
The terms with $j=k$ vanish and can be restored. Expanding the square
gives the two terms $q_2\otimes q_0$ and $q_0\otimes q_2$,
minus $2q_1\otimes q_1$, all tested against $\chi(x-x')$.
The last term is zero. Evenness of $\chi$ makes the other two equal.
Their positivity gives the upper bound.
\end{proof}

If $\chi(0)=1$ and there are $m$ copies at each of $x_*\pm\ii d$,
their two ordered cross populations contribute
$8m^2w((x_*-X)/h)^2d^2$ to $S_2$.
They contribute $2mw((x_*-X)/h)d^2$ to $Q_w$.
Thus the test detects a transverse pair, but its exact reduction
prevents interpreting it as independently acquired two-point
information.

The same mechanism extends to all even difference powers. For $k\ge1$,
\begin{equation}\label{eq:higher-moments}
 S_{2k}(\chi)=
 \sum_{r=0}^k \binom{2k}{2r}
       \iint\chi(x-x')\dd q_{2k-2r}(x)\dd q_{2r}(x').
\end{equation}
This follows by binomial expansion; all odd factors vanish.
Since $|y_j-y_k|\le2$,
$S_{2k}(\chi)\le2^{2k-2}S_2(\chi)$.
These are identities and bounds for this particular family of tests;
they do not reconstruct the full marked pair measure.

\subsection{A configuration generating function with an order bound}

Put $a_j=w_jy_j^2\in[0,1]$ and
\begin{equation}\label{eq:configuration-product}
 Z(\eta)=\prod_{j=1}^N(1+\eta a_j)
        =\sum_{k=0}^N e_k\eta^k,\qquad
 p_r=\sum_j a_j^r,\qquad p_1=Q_w.
\end{equation}
The $e_k$ are elementary symmetric polynomials of the $a_j$.
Differentiating the finite product gives, first near $\eta=0$,
\[
 \frac{Z'(\eta)}{Z(\eta)}
 =\sum_j\frac{a_j}{1+\eta a_j}
 =\sum_{r\ge1}(-1)^{r-1}p_r\eta^{r-1}.
\]
Comparing coefficients proves Newton's recurrence
\begin{equation}\label{eq:newton}
 e_0=1,\qquad
 ke_k=\sum_{r=1}^k(-1)^{r-1}p_re_{k-r}.
\end{equation}
It is exact finite algebra; no closure assumption enters.

\begin{proposition}[Product comparison and growing-order truncation]\label{prop:product}
For $\eta\ge0$,
\begin{equation}\label{eq:product-bounds}
 \log(1+\eta)Q_w\le\log Z(\eta)\le\eta Q_w .
\end{equation}
For every integer $K\ge0$,
\begin{equation}\label{eq:order-tail}
 0\le Z(\eta)-\sum_{k=0}^Ke_k\eta^k
 \le \ee^{\eta Q_w}
       \frac{(\eta Q_w)^{K+1}}{(K+1)!},
\end{equation}
where $e_k=0$ for $k>N$. If $Q_w\le B$ and $0<\varepsilon<1$,
the right side is at most $\varepsilon$ whenever $k=K+1$ satisfies
\begin{equation}\label{eq:order-choice}
 k\ge\max\left\{1,\ 2\ee\eta B,\
           \frac{\eta B+\log(1/\varepsilon)}{\log2}\right\}.
\end{equation}
\end{proposition}
\begin{proof}
On $a\in[0,1]$, concavity of $\log(1+\eta a)$ puts its graph
above the chord $a\log(1+\eta)$ and below its tangent $\eta a$.
Summing proves \eqref{eq:product-bounds}.
Expansion of $(\sum a_j)^k$ contains every product with distinct labels
$k!$ times and additional nonnegative terms, so $e_k\le Q_w^k/k!$.
For $t\ge0$,
$\sum_{j=k}^\infty t^j/j!\le t^k\ee^t/k!$, since
$(k+\ell)!\ge k!\ell!$. This proves \eqref{eq:order-tail}.
Finally $k!\ge(k/\ee)^k$ makes its right side at most
$\ee^{\eta B}2^{-k}$ under the first two size conditions, and the
last condition gives $\varepsilon$.
\end{proof}

For complex $\eta$, use the logarithm branch equal to zero at $\eta=0$.
If $q=|\eta|\max_j a_j<1$, the same finite product gives
\begin{equation}\label{eq:complex-log-tail}
 \left|\log Z(\eta)-
       \sum_{r=1}^K\frac{(-1)^{r-1}}r\,p_r\eta^r\right|
 \le
 \frac{|\eta|^{K+1}p_{K+1}}{(K+1)(1-q)} .
\end{equation}
For each factor, expand $\log(1+\eta a_j)$, use
$1/r\le1/(K+1)$ in the tail, and sum the geometric series.
Thus increasing the factorial order is permitted only with a
simultaneous radius or tail payment, not from fixed-order identities
alone.

A local count can supply $B$. For example, if the population in
every relevant interval of length at most one is bounded by
$A_{\mathrm{count}}\log(3+|X|+h)$, and
$\supp w\subset[-1,1]$, covering $[X-h,X+h]$ gives
\[
 Q_w\le N_w\le
 A_{\mathrm{count}}(\lceil2h\rceil+1)\log(3+|X|+h).
\]
This is a conditional count corollary, with its constant retained.
The canonical disk estimate in Section~\ref{sec:canonical} supplies a
different explicit population bound in its stated regime.

\subsection{Defect signals and the cost of averaging}

On a plateau where $w=1$, a simple pair $x_*\pm\ii d$ contributes
\begin{equation}\label{eq:scalar-signals}
 Q_w\ge2d^2,\qquad
 \log Z(\eta)\ge2\log(1+\eta d^2).
\end{equation}
All other terms are nonnegative. Both observables are zero for every
all-real configuration, including real multiple zeros. They test
nonreality, not simplicity.

If the center $X$ is averaged uniformly on an interval of length $L$
which contains the full translate of the window around $x_*$, the
pair's respective contributions become
\begin{align}
 &\frac{2d^2h}{L}\int_\R w(t)\dd t,\qquad
 \frac{2h}{L}\int_\R\log(1+\eta d^2w(t))\dd t,\label{eq:average-scalar}\\
 &\frac{8d^2h}{L}\int_\R w(t)^2\dd t
       \quad\text{for its mutual contribution to }S_2
       \text{ when }\chi(0)=1 .\notag
\end{align}
These follow by the substitution $t=(x_*-X)/h$.
Averaged and every-center statements therefore have different
single-pair margins. With complete height and transverse-scale
coverage, vanishing of these nonnegative tests is simply a
reformulation of the assertion that all canonical zeros are real,
equivalently RH. That reformulation supplies no new estimate and is
not supported by sampling checks.
